\documentclass[11pt]{article}
\usepackage[T1]{fontenc}
\usepackage{lmodern}
\usepackage{microtype}
\usepackage[a4paper,margin=28mm]{geometry}
\usepackage{amsmath,amssymb,amsthm,mathtools}
\usepackage{graphicx}
\usepackage{xcolor}
\usepackage[hidelinks]{hyperref}
\usepackage{enumitem}

\definecolor{packetblue}{RGB}{31,78,121}
\hypersetup{colorlinks=true,linkcolor=packetblue,urlcolor=packetblue,citecolor=packetblue}
\setlist{nosep}
\newtheorem{theorem}{Theorem}
\newtheorem{lemma}[theorem]{Lemma}
\newtheorem{remark}[theorem]{Remark}
\newcommand{\Exh}{\operatorname{Exh}}
\newcommand{\Fin}{\operatorname{Fin}}
\newcommand{\N}{\mathbb N}

\title{Every Density-Like Ideal Has a\\Strongly-Density-Like Representation}
\author{Candidate solution packet for Question 5 of Kwela--Leonetti}
\date{9 August 2026}

\begin{document}
\maketitle

\begin{center}
\fcolorbox{packetblue}{packetblue!6}{\parbox{0.88\textwidth}{
\textbf{Status.} Candidate full solution, likely valid; not human verified. The result gives a negative answer to the source question by proving a stronger representation theorem.}}
\end{center}

\begin{abstract}
Kwela and Leonetti ask whether some density-like ideal fails to admit any strongly-density-like lower semicontinuous submeasure representation. We prove that no such ideal exists. Given an arbitrary density-like lower semicontinuous submeasure $\varphi$, we select scales adapted to its density modulus and interpolate them by a concave gauge $f$. The composition $f\circ\varphi$ has a fixed linear density modulus, while its exhaustive ideal is unchanged.
\end{abstract}

\section{The source question}

The source is Adam Kwela and Paolo Leonetti, \emph{Density-Like and Generalized Density Ideals}, arXiv:2001.00223, published in the \emph{Journal of Symbolic Logic} 87 (2022), 228--251. Section 6 concludes with the following question (Question 5, source PDF p.~25).

\begin{quote}
Does there exist a density-like ideal $\mathcal I$ such that $\mathcal I\ne\Exh(\varphi)$ for each strongly-density-like lscsm $\varphi$?
\end{quote}

The complete source statement and its immediate context are reproduced on the next page.

\clearpage
\begin{center}
{\large\bfseries Source statement (paper p.~25)}

\vspace{12mm}
\includegraphics[width=\textwidth]{figures/open_problem_crop.png}
\end{center}

\vspace{8mm}
\subsection*{Source reference}

A. Kwela and P. Leonetti, \emph{Density-Like and Generalized Density Ideals}, Journal of Symbolic Logic \textbf{87} (2022), 228--251; arXiv:2001.00223; doi:10.1017/jsl.2021.95.
\clearpage
\section{Definitions and result}

Write $\omega=\{0,1,2,\ldots\}$. A lower semicontinuous submeasure (lscsm) is a monotone, subadditive map
\[
 \varphi:\mathcal P(\omega)\longrightarrow[0,\infty]
\]
such that $\varphi(\varnothing)=0$, $\varphi(\{n\})<\infty$ for every $n$, and
\[
 \varphi(A)=\lim_{m\to\infty}\varphi(A\cap m).
\]
Its exhaustive ideal is
\[
 \Exh(\varphi)=\left\{A\subseteq\omega:
 \|A\|_\varphi:=\lim_{m\to\infty}\varphi(A\setminus m)=0\right\}.
\]

The lscsm $\varphi$ is \emph{density-like} if, for every $\varepsilon>0$, there is $\delta>0$ such that every sequence $(F_k)$ of pairwise disjoint finite sets satisfying $\varphi(F_k)<\delta$ has an infinite $I\subseteq\omega$ for which
\[
 \varphi\left(\bigcup_{k\in I}F_k\right)<\varepsilon.
\]
It is \emph{strongly-density-like} if one may always take $\delta=c\varepsilon$ for a fixed constant $c>0$ independent of $\varepsilon$.

\subsection*{Idea of the proof}

The original density-like modulus may be very nonlinear, but only countably many target scales are needed. We choose rapidly decreasing scales $t_n$ so that input size below $t_{n+1}$ activates the density conclusion at output target $t_n$. We then change coordinates by a piecewise-affine gauge satisfying $f(t_n)=2^{-n}$. The scale separation makes $f$ concave, hence subadditive, while the dyadic values turn the arbitrary original modulus into the fixed estimate $\delta=\varepsilon/4$. Finally, continuity at zero and positivity away from zero ensure that this change of coordinates preserves the exhaustive ideal exactly.

\begin{theorem}[Concave-gauge uniformization]\label{thm:main}
Let $\varphi$ be a density-like lscsm on $\omega$. There exists a bounded, continuous, increasing, concave function
\[
 f:[0,\infty]\longrightarrow[0,1]
\]
with $f(0)=0$ and $f(t)>0$ for every $t>0$, such that
\[
 \psi=f\circ\varphi
\]
is a strongly-density-like lscsm with witnessing constant $c=1/4$, and
\[
 \Exh(\psi)=\Exh(\varphi).
\]
Consequently, every density-like ideal has a strongly-density-like lscsm representation, so the answer to Question 5 is \emph{no}.
\end{theorem}

\section{Construction of the gauge}

Set $t_0=1$. Suppose $t_n>0$ has been selected. By density-likeness, choose $d_n>0$ such that every pairwise disjoint sequence $(F_k)$ of finite sets with
\[
 \varphi(F_k)<d_n\qquad(k\in\omega)
\]
admits an infinite $I\subseteq\omega$ satisfying
\[
 \varphi\left(\bigcup_{k\in I}F_k\right)<t_n.
\]
Now choose
\begin{equation}\label{eq:scales}
 0<t_{n+1}<\min\{d_n,t_n/4\}.
\end{equation}
Then $t_n\downarrow0$.

Define $f(0)=0$, set
\[
 f(t_n)=2^{-n}\qquad(n\ge0),
\]
make $f$ affine on every interval $[t_{n+1},t_n]$, and set $f(t)=1$ for $t\ge t_0$ and $f(\infty)=1$. This gives a continuous, strictly increasing function on $[0,t_0]$ that is constant thereafter and vanishes only at zero.

\begin{lemma}\label{lem:concave}
The function $f$ is concave and subadditive on $[0,\infty]$.
\end{lemma}

\begin{proof}
The slope on $[t_{n+1},t_n]$ is
\[
 s_n=\frac{2^{-(n+1)}}{t_n-t_{n+1}}.
\]
By \eqref{eq:scales},
\[
 \frac{s_{n+1}}{s_n}
 =\frac12\frac{t_n-t_{n+1}}{t_{n+1}-t_{n+2}}
 >\frac12\frac{3t_n/4}{t_n/4}
 =\frac32.
\]
Thus the slopes increase as the intervals approach the origin, or equivalently decrease from left to right. The slope of the constant part $[t_0,\infty)$ is zero. Hence $f$ is concave.

For finite $x,y\ge0$ with $x+y>0$, concavity and $f(0)=0$ give
\[
 f(x)\ge\frac{x}{x+y}f(x+y),
 \qquad
 f(y)\ge\frac{y}{x+y}f(x+y).
\]
Adding yields $f(x+y)\le f(x)+f(y)$. The cases involving infinity are immediate from $f(\infty)=1$ and $0\le f\le1$.
\end{proof}

\section{Proof of the theorem}

By monotonicity and subadditivity of $\varphi$ and $f$, the composition $\psi=f\circ\varphi$ is monotone, vanishes at the empty set, and is subadditive. Moreover,
\[
 \psi(\{n\})=f(\varphi(\{n\}))\le1<\infty.
\]
For $A\subseteq\omega$, lower semicontinuity of $\varphi$ and continuity of $f$ imply
\[
 \psi(A)=f(\varphi(A))
 =f\left(\lim_m\varphi(A\cap m)\right)
 =\lim_m f(\varphi(A\cap m))
 =\lim_m\psi(A\cap m).
\]
This also covers $\varphi(A)=\infty$, because $f$ is already equal to $1$ on $[t_0,\infty]$. Hence $\psi$ is a lscsm.

We next compare the exhaustive ideals. The tail values $\varphi(A\setminus m)$ decrease to $\|A\|_\varphi$, possibly infinity. Continuity of $f$ gives
\begin{equation}\label{eq:tail}
 \|A\|_\psi
 =\lim_m f(\varphi(A\setminus m))
 =f(\|A\|_\varphi).
\end{equation}
Since $f(u)=0$ exactly when $u=0$, equation \eqref{eq:tail} proves
\[
 \Exh(\psi)=\Exh(\varphi).
\]

It remains to prove the strong-density estimate. Fix $\varepsilon>0$ and a pairwise disjoint sequence $(F_k)$ of finite sets such that
\begin{equation}\label{eq:stronghyp}
 \psi(F_k)<\varepsilon/4\qquad(k\in\omega).
\end{equation}
If $\varepsilon>1$, then every infinite $I\subseteq\omega$ works, because $\psi\le1<\varepsilon$.

Suppose $0<\varepsilon\le1$. Choose the unique $n\ge0$ such that
\begin{equation}\label{eq:dyadic}
 2^{-n}\le\varepsilon<2^{1-n}.
\end{equation}
By \eqref{eq:stronghyp} and \eqref{eq:dyadic},
\[
 \psi(F_k)<\frac{\varepsilon}{4}
 <2^{-n-1}=f(t_{n+1}).
\]
As $f$ is increasing and strictly increasing below $t_0$, this implies
\[
 \varphi(F_k)<t_{n+1}<d_n
 \qquad(k\in\omega).
\]
The choice of $d_n$ therefore supplies an infinite $I\subseteq\omega$ with
\[
 \varphi\left(\bigcup_{k\in I}F_k\right)<t_n.
\]
Applying $f$ and using \eqref{eq:dyadic},
\[
 \psi\left(\bigcup_{k\in I}F_k\right)
 <f(t_n)=2^{-n}\le\varepsilon.
\]
Thus $c=1/4$ witnesses that $\psi$ is strongly-density-like. Together with \eqref{eq:tail}, this completes the proof of Theorem~\ref{thm:main}. \qed

\section{Checks, scope, and novelty screen}

The proof is representation-theoretic: it does not assert that a given density-like representation is already strongly-density-like. Indeed, the source paper exhibits density-like lscsms that are not strongly-density-like. The theorem says that the same exhaustive ideal always has another representation that is strongly-density-like.

The numerical script \texttt{code/check\_concave\_gauge.py} checks a finite model with 41 rapidly separated scales. It verifies all adjacent slope inequalities, 50,000 random concavity inequalities, 50,000 random subadditivity inequalities, and the dyadic threshold identities at 31 scales. These checks are diagnostic only; the proof above is analytic.

A bounded literature search performed on 9 August 2026 used the exact question wording, the term ``strongly-density-like,'' the paper title, arXiv, and OpenAlex. It located the source paper but no later mathematical answer. OpenAlex's only listed citing record was cover/back matter for the same journal issue. This supports review of the result as a potentially new answer, but it is not an exhaustive priority claim.

\medskip
\noindent\textbf{Human review recommendation.} Check especially the direction of the slope comparison in Lemma~2, the extended-real continuity step when $\varphi(A)=\infty$, and the strict dyadic inequalities leading to the constant $1/4$. A broader MathSciNet or Zentralblatt search is recommended before making any priority claim.

\end{document}
